
\medskip
\noindent\textbf{0.3.5.}
\phantomsection
\label{sga3:VIIB:0.3.5}
Consider three pseudocompact \(A\)-modules \(L,M,N\), where \(N\) has
finite length. Give the \(A\)-module
\(\operatorname{Hom}_c(M,N)\) the discrete topology. Every element
\[
  \psi\in
  \operatorname{Hom}_c
    \bigl(L,\operatorname{Hom}_c(M,N)\bigr)
\]
defines a continuous bilinear map
\[
  \psi':L\times M\longrightarrow N,
  \qquad
  (\ell,m)\longmapsto\psi(\ell)(m).
\]
This gives a natural isomorphism
\begin{equation}
  \operatorname{Hom}_c
    \bigl(L,\operatorname{Hom}_c(M,N)\bigr)
  \xrightarrow{\sim}
  \operatorname{Hom}_c
    \bigl(L\widehat{\otimes}_A M,N\bigr).
  \tag{1}
  \label{sga3:VIIB:0.3.5-eq1}
\end{equation}
This is another characterization of
\(\operatorname{Hom}_c(L\widehat{\otimes}_A M,N)\), which we now use
when \(M\) is the projective limit of a filtered projective system
\((M_i)\) of pseudocompact \(A\)-modules. By (1) and \SGARefLink{sga3:VIIB:0.2.F}{0.2.F}~\SGARefLink{sga3:vii:paragraph:5-1:equation:ii}{(ii)}, there
are natural isomorphisms
\begin{align}
  \operatorname{Hom}_c
    \left(L\widehat{\otimes}_A\varprojlim_iM_i,N\right)
  &\simeq
  \operatorname{Hom}_c
    \left(L,\operatorname{Hom}_c(\varprojlim_iM_i,N)\right)
  \notag\\
  &\simeq
  \operatorname{Hom}_c
    \left(L,\varinjlim_i\operatorname{Hom}_c(M_i,N)\right).
  \tag{2}
  \label{sga3:VIIB:0.3.5-eq2}
\end{align}

Moreover, because
\(\varinjlim_i\operatorname{Hom}_c(M_i,N)\) is discrete, every
continuous map from \(L\) into this module factors through a
finite-length quotient of \(L\). Hence the natural map
\begin{equation}
  \varinjlim_i
  \operatorname{Hom}_c
    \bigl(L,\operatorname{Hom}_c(M_i,N)\bigr)
  \xrightarrow{\sim}
  \operatorname{Hom}_c
    \left(L,\varinjlim_i\operatorname{Hom}_c(M_i,N)\right)
  \tag{3}
  \label{sga3:VIIB:0.3.5-eq3}
\end{equation}
is an isomorphism. Finally, (1) and \SGARefLink{sga3:VIIB:0.2.F}{0.2.F}~\SGARefLink{sga3:vii:paragraph:5-1:equation:ii}{(ii)} again give natural
isomorphisms
\begin{align}
  \varinjlim_i
  \operatorname{Hom}_c
    \bigl(L,\operatorname{Hom}_c(M_i,N)\bigr)
  &\simeq
  \varinjlim_i
  \operatorname{Hom}_c
    \bigl(L\widehat{\otimes}_A M_i,N\bigr)
  \notag\\
  &\simeq
  \operatorname{Hom}_c
    \left(\varprojlim_i
      (L\widehat{\otimes}_A M_i),N\right).
  \tag{4}
  \label{sga3:VIIB:0.3.5-eq4}
\end{align}
Combining \SGARefLink{sga3:VIIB:0.3.5-eq2}{(2)}, \SGARefLink{sga3:VIIB:0.3.5-eq3}{(3)}, and \SGARefLink{sga3:VIIB:0.3.5-eq4}{(4)} yields the following result.

\medskip
\noindent\textbf{Proposition.}
\phantomsection
\label{sga3:VIIB:0.3.5-proposition}
{\itshape Let \((M_i)\) be a filtered projective system of objects of
\(\operatorname{PC}(A)\), let
\(L\in\operatorname{PC}(A)\), and let
\(N\in\operatorname{LF}(A)\). There is an isomorphism, functorial in
\(N\),
\[
  \operatorname{Hom}_c
    \left(L\widehat{\otimes}_A\varprojlim_iM_i,N\right)
  \simeq
  \operatorname{Hom}_c
    \left(\varprojlim_i(L\widehat{\otimes}_A M_i),N\right),
\]
and consequently an isomorphism
\[
  L\widehat{\otimes}_A\varprojlim_iM_i
  \simeq
  \varprojlim_i(L\widehat{\otimes}_A M_i).
\]
Thus the completed tensor product commutes with filtered projective
limits.}

\medskip
\noindent\textbf{0.3.6.}
\phantomsection
\label{sga3:VIIB:0.3.6}
In particular, the completed tensor
product commutes with arbitrary products. For example, since \(A\) is
the product of its local components \(A_{\mathfrak m}\) (\SGARefLink{sga3:VIIB:0.1.1}{0.1.1}), every
pseudocompact \(A\)-module
\[
  M\simeq A\widehat{\otimes}_A M
\]
identifies with the product of the modules
\[
  M_{\mathfrak m}
  =
  A_{\mathfrak m}\widehat{\otimes}_A M,
\]
which are called the local components of \(M\).

Likewise, let \(M\) and \(N\) be pseudocompact \(A\)-modules. Recall
from \SGARefLink{sga3:VIIB:0.2.2}{0.2.2} that
\[
  M^\dagger=\operatorname{Hom}_c(M,A).
\]
Consider the map
\[
  \varphi:
  M^\dagger\otimes_A N^\dagger
  \longrightarrow
  (M\widehat{\otimes}_A N)^\dagger
\]
which sends \(f\otimes g\) to the map
\[
  m\widehat{\otimes}n\longmapsto f(m)g(n)
\]
from \(M\widehat{\otimes}_A N\) to \(A\). The map \(\varphi\) is
bijective when \(M\simeq A\).

\medskip
\noindent\textbf{Corollary.}
\phantomsection
\label{sga3:VIIB:0.3.6-corollary}
\textit{If \(A\) is Artinian, then \(\varphi\) is an isomorphism
whenever \(M\) is topologically free, or more generally projective.}

Indeed, with \(N\) fixed, each of the functors
\[
  M\longmapsto(M\widehat{\otimes}_A N)^\dagger
  \qquad\text{and}\qquad
  M\longmapsto M^\dagger\otimes_A N^\dagger
\]
takes arbitrary products to direct sums, by the preceding result and
\SGARefLink{sga3:VIIB:0.2.F}{0.2.F}.

\medskip
\noindent\textbf{Remark 0.3.6.A.}\footnote{The editors added this
remark, which will be useful in \SGARefLink{sga3:VIIB:2.3.1}{2.3.1}.}
\phantomsection
\label{sga3:VIIB:0.3.6.A}
\textit{The same use of \SGARefLink{sga3:VIIB:0.2.F}{0.2.F} gives the following result. Let \(A\) be
an Artinian ring, let \(M,Q\in\operatorname{PC}(A)\), and let
\(N\in\operatorname{LF}(A)\). If \(Q\) is projective, then there are
natural isomorphisms}
\[
  \operatorname{Hom}_c(M,Q)
  \xrightarrow{\sim}
  \operatorname{Hom}_A(Q^\dagger,M^\dagger),
  \qquad
  Q^\dagger\otimes_A N
  \xrightarrow{\sim}
  \operatorname{Hom}_c(Q,N).
\]

\medskip
\noindent\textbf{0.3.7.}
\phantomsection
\label{sga3:VIIB:0.3.7}
For every \(\mathfrak m\in\Upsilon(A)\), the functor
\[
  M\longmapsto A_{\mathfrak m}\widehat{\otimes}_A M
\]
is plainly exact. Since every projective pseudocompact \(A\)-module
\(P\) is a product of modules of the form \(A_{\mathfrak m}\), it
follows that
\[
  M\longmapsto P\widehat{\otimes}_A M
\]
is exact whenever \(P\) is projective. The converse is also true.

\medskip
\noindent\textbf{Proposition.}
\phantomsection
\label{sga3:VIIB:0.3.7-proposition}
{\itshape Let \(A\) be a pseudocompact ring and \(P\) a pseudocompact
\(A\)-module. The following conditions are equivalent:
\begin{enumerate}
  \item[\textup{(i)}]
  \(P\) is a projective object of \(\operatorname{PC}(A)\);
  \item[\textup{(ii)}]
  every local component \(P_{\mathfrak m}\) is a topologically free
  \(A_{\mathfrak m}\)-module;
  \item[\textup{(iii)}]
  the functor \(M\mapsto P\widehat{\otimes}_A M\) is exact.
\end{enumerate}}

The equivalence of \SGARefLink{sga3:vii:paragraph:5-1:equation:i}{(i)} and \SGARefLink{sga3:vii:paragraph:5-1:equation:ii}{(ii)} follows from \SGARefLink{sga3:VIIB:0.2.F}{0.2.F}~\SGARefLink{sga3:VIIB:0.2.F:item:iii}{(iii)} and \SGARefLink{sga3:VIIB:0.2.1}{0.2.1},
and we have just proved \SGARefLink{sga3:VIIB:0.2.1:item:ii}{(ii)}~\(\Rightarrow\)~\SGARefLink{sga3:VIIB:0.2.F:item:iii}{(iii)}. Suppose that the
functor \(M\mapsto P\widehat{\otimes}_A M\) is exact. Since
\(P\widehat{\otimes}_A M\) is the product of its local components,
\[
  (P\widehat{\otimes}_A M)_{\mathfrak m}
  \simeq
  P_{\mathfrak m}
    \widehat{\otimes}_{A_{\mathfrak m}}M_{\mathfrak m},
\]
we reduce to the case in which \(A\) is local. We shall prove that
\(P\) is then topologically free.

Let \(\mathfrak m\) be the maximal ideal of \(A\). The quotient
\(P/\overline{\mathfrak m}P\) is a linearly compact vector space over
\(A/\mathfrak m\), and hence a product of copies of \(A/\mathfrak m\)
(\SGARefLink{sga3:VIIB:0.2.2}{0.2.2}). Therefore there is a family \((A_i)_{i\in I}\) of copies of
\(A\) and an isomorphism
\[
  \phi:
  \prod_{i\in I}(A_i/\mathfrak m)
  \xrightarrow{\sim}
  P/\overline{\mathfrak m}P.
\]
Since \(\prod_{i\in I}A_i\) is projective, there is a commutative
square
\begin{center}
  \phantomsection
  \label{sga3:VIIB:0.3.7-projective-lift-square}
  \begin{tikzcd}[row sep=3em,column sep=4.8em]
\displaystyle\prod_{i\in I}A_i \arrow[r,"\psi"] \arrow[d,"p"']
  & P \arrow[d,"q"] \\
\displaystyle\prod_{i\in I}(A_i/\mathfrak m)
  \arrow[r,"\varphi","{\sim}"']
  & P/\mathfrak mP
\end{tikzcd}

\end{center}
where \(p\) and \(q\) are the canonical projections. Apply Nakayama's
lemma to \(\operatorname{Coker}\psi\). Since
\[
  (A/\mathfrak m)\widehat{\otimes}_A\psi
\]
is precisely \(\phi\), the map \(\psi\) is
surjective.

Put
\[
  B=\prod_{i\in I}A_i,
  \qquad
  N=\ker\psi.
\]
There is then the following commutative exact diagram:
\begin{center}
  \phantomsection
  \label{sga3:VIIB:0.3.7-nakayama-exact-diagram}
  \begin{tikzcd}[row sep=2.05em,column sep=1.65em]
{} & {} & {} &
  0
    \arrow[d]
    \arrow[ddlll,dashed,rounded corners=10pt,
      to path={
        -- ([xshift=6em]\tikztostart.east)
        |- ([yshift=.8em]\tikztotarget.north)
        -- (\tikztotarget)
      }]
  & {} \\
{} & \mathfrak m\widehat\otimes_A N \arrow[r] \arrow[d]
  & \mathfrak m\widehat\otimes_A B \arrow[r] \arrow[d]
  & \mathfrak m\widehat\otimes_A P \arrow[r] \arrow[d]
  & 0 \\
0
  \arrow[r]
  \arrow[dr,dashed,out=-90,in=180,looseness=1.35]
  & A\widehat\otimes_A N \arrow[r] \arrow[d]
  & A\widehat\otimes_A B \arrow[r,"\psi"] \arrow[d]
  & A\widehat\otimes_A P \arrow[r] \arrow[d]
  & 0 \\
{} & (A/\mathfrak m)\widehat\otimes_A N \arrow[r]
  & (A/\mathfrak m)\widehat\otimes_A B \arrow[r,"\varphi"]
  & (A/\mathfrak m)\widehat\otimes_A P \arrow[r]
  & 0 \qquad .
\end{tikzcd}

\end{center}
The snake lemma, applied to the first two rows, shows that in the
bottom row the morphism
\[
  (A/\mathfrak m)\widehat{\otimes}_A N
  \longrightarrow
  (A/\mathfrak m)\widehat{\otimes}_A B
\]
is a monomorphism. But \(\phi\) is an isomorphism, so
\((A/\mathfrak m)\widehat{\otimes}_A N=0\). Hence \(N=0\) by \SGARefLink{sga3:VIIB:0.3.3}{0.3.3},
and \(\psi\) is an isomorphism.

\medskip
\noindent\textbf{0.3.8. Corollary.}
\phantomsection
\label{sga3:VIIB:0.3.8}
\textit{Let \(A\) be a complete Noetherian local ring and \(P\) a
pseudocompact \(A\)-module. Then \(P\) is topologically free if and
only if \(P\) is flat over \(A\).}

Indeed, the canonical map
\[
  M\otimes_A P\longrightarrow M\widehat{\otimes}_A P
\]
is bijective when \(M=A\), and hence also when \(M\) is Noetherian:
take a finite presentation of \(M\) and use the right exactness of
both the tensor product and the completed tensor product.

Now \(P\) is flat if and only if
\[
  M\longmapsto M\otimes_A P
\]
is exact on Noetherian modules. Likewise, the proof of \SGARefLink{sga3:VIIB:0.3.7}{Proposition
0.3.7} shows that \(P\) is topologically free provided the sequence
\[
  0\longrightarrow
  \mathfrak m\widehat{\otimes}_A P
  \longrightarrow
  A\widehat{\otimes}_A P
  \longrightarrow
  (A/\mathfrak m)\widehat{\otimes}_A P
  \longrightarrow0
\]
is exact. Thus \(P\) is topologically free if and only if the functor
\[
  M\longmapsto M\widehat{\otimes}_A P
\]
is exact on Noetherian modules. The corollary follows from the equality
\[
  M\otimes_A P=M\widehat{\otimes}_A P
\]
established above.

